\documentclass[10pt,a4paper]{amsart}
\usepackage[margin=19mm]{geometry}
\usepackage{amsmath,amssymb,mathtools}
\usepackage[scr=rsfs,cal=euler]{mathalfa}
\usepackage{xfrac}
\usepackage[unicode,hidelinks,bookmarks=false]{hyperref}
\newcommand{\ie}{i.e.\ }
\newcommand{\cf}{cf.\ }
\newcommand{\resp}{respectively\ }
\newcommand{\Subsection}{\subsection}
\providecommand{\nobreakdash}{\nobreak\hbox{-}}
\setlength{\emergencystretch}{2em}
\multlinegap0pt
\usepackage{amssymb}
\usepackage{dsfont}
\usepackage[shortlabels]{enumitem}
\newtheorem{lem}{Lemma}
\title{A fixed point theorem for the action of linear higher rank algebraic groups over local fields on symmetric spaces of infinite dimension and finite rank}
\author{Federico Viola}
\date{Source: arXiv:2505.05220v2 (CC BY 4.0)}
\begin{document}
\maketitle

\noindent\emph{Source and licence.} A fixed point theorem for the action of linear higher rank algebraic groups over local fields on symmetric spaces of infinite dimension and finite rank. Version arXiv:2505.05220v2 (CC BY 4.0), distributed by arXiv under the Creative Commons Attribution 4.0 International licence, http://creativecommons.org/licenses/by/4.0/, as recorded in the arXiv licence metadata for this exact version and shown on the abstract page https://arxiv.org/abs/2505.05220v2. Full licence notice: \texttt{licenses/arxiv-2505.05220-CC-BY-4.0.txt}.\par\smallskip

\noindent\textit{Editorial scope.} Counted unit: the article's lem soluble (source0.tex lines 343--345). The article's complete original proof of it is delivered with it (source0.tex lines 347--358, one \texttt{proof} environment of 12 lines). No mathematics is written, repaired or re-derived: the statement and the proof are the article's own lines, in the article's own notation, and no retained line is substituted or re-flowed.  No further source-local context is needed: every cross-reference of every retained line resolves inside the two delivered spans.  Nothing was omitted from any retained span, so the package carries no omission record.  No retained line contains a citation, so the wrapper carries no bibliography and no bibliography entry.  No retained line contains a figure environment, an image-inclusion command or drawing code, so no image is delivered and the package declares no asset dependency.  Editorial formatting only, in addition: an AMS \texttt{amsart} preamble that restates the article's own declarations and macros used by the retained lines, namely the environments \texttt{lem} on separate counters, together with the standard packages those lines need.  The retained environments print in this document as Lemma 1 in order of appearance, whereas the article prints the same items with the numbering of its own full document.  The complete native source of the article as distributed in the arXiv e-print for this exact version is bundled unchanged as \texttt{sources/177800/source0.tex}.
\begin{lem} \label{soluble}
    Let $\Gamma$ be a locally compact group with Property (FH), and let $H=N\rtimes_\phi Q$ be a semidirect product of topological groups. Assume that $N$ is soluble and that every Abelian quotient that arises from its derived series is isomorphic (as a topological group) to the additive group of a Hilbert space, where the action of $Q$ (induced from $\phi$) preserves a scalar product. Then any continuous homomorphism $f:\Gamma\rightarrow H$ has image contained in a conjugate of $Q$.
\end{lem}

\begin{proof}
    We prove the lemma by induction on the solubility length of $N$. The base case will be the trivial case with solubility length zero, where $N$ is the trivial group, and $H=Q$. For the inductive step, we assume that $N$ has solubility length $n\geq 1$ and that the lemma is true in all cases with solubility length $\leq n-1$. \\

    Let $N':=[N,N]$ be the commutator subgroup of $N$, which has solubility length $n-1$, let $A=N/N'$ be the abelianization of $N$, and let $\mathrm{ab}: N\rightarrow A$ be the natural projection (with kernel $N'$). We recall that $Q$ acts on $N$ by the map $\phi: Q\rightarrow\mathrm{Aut}(N)$ that appears in the semidirect product and that associates to every $q\in Q$ the conjugation by $q$. Since the commutator subgroup $N'$ is characteristic, it is preserved by the action of $Q$, which means that the action descends to the quotient $A$. We call $\phi_{\mathrm{ab}}$ this induced action. By hypothesis, $A$ is isomorphic (as a topological group) to the additive group of a Hilbert space, and the action $\phi_{\mathrm{ab}}$ of $Q$ preserves a scalar product in this Hilbert space. \\

    We can now define an action of $\Gamma$ on $A$ as follows: for every $g\in \Gamma$ there exist unique $n(g)\in N, q(g)\in Q$ such that $f(g)=n(g)q(g)$; then we set $$F(g)(a)=\mathrm{ab}(n(g))\phi_{\mathrm{ab}}(q(g))(a)\ \ \ \forall a\in A.$$
    
    The action of $Q$ is by automorphisms, so by linear operators, while the action of $N$ is by translations. Therefore, the action of $\Gamma$ is by affine transformations, and as the action of $Q$ preserves a scalar product, the whole action of $\Gamma$ is by isometries with respect to the metric induced by this scalar product on $A$. \\
    
    Using Property (FH), we find that the action of $\Gamma$ on $A$ fixes a point $a_0$, which means that if we conjugate $f$ by an element $n_0\in N$ with $\mathrm{ab}(n_0)=a_0$ we get an action that fixes the neutral element of $A$, which means that a conjugate of the homomorphism $f$ has image contained in $N'\rtimes Q$. Now, since $N'$ has solubility length $n-1$, by inductive hypothesis this conjugate has image contained in a conjugate of $Q$, so in the end we find that $f$ has image contained in a conjugate of $Q$. \\
    
\end{proof}

\end{document}
